% Generated draft: requires mathematical and source review.
\documentclass[11pt]{article}
\usepackage[margin=25mm]{geometry}
\usepackage{fontspec}
\setmainfont{Latin Modern Roman}
\usepackage{amsmath,amssymb,mathtools,unicode-math}
\setmathfont{Latin Modern Math}
\usepackage{enumitem,needspace,xurl,hyperref}
\setlength{\parindent}{0pt}
\setlength{\parskip}{5pt}
\setlength{\emergencystretch}{3em}
\newcommand{\N}{\mathbb{N}}
\newcommand{\Z}{\mathbb{Z}}
\newcommand{\Q}{\mathbb{Q}}
\newcommand{\R}{\mathbb{R}}
\newcommand{\C}{\mathbb{C}}
\newcommand{\F}{\mathbb{F}}
\newcommand{\A}{\mathbb{A}}
\newcommand{\PP}{\mathbb{P}}
\newcommand{\E}{\mathbb{E}}
\newcommand{\eps}{\varepsilon}
\newcommand{\dd}{\,\mathrm{d}}
\begin{document}
\section*{UnsolvedMath ID: 2586}
\textbf{Source problem code:} catalogue entry
\subsection{Definitions and mathematical statement}
Take the nontrivial-degree interpretation $n\geq2$. Let $d\in\Z$ satisfy $p^n\nmid d$ for every prime $p$, and suppose $x^n-d$ is irreducible over $\Q$. Choose a root $\theta$ with $\theta^n=d$ and put $K=\Q(\theta)$, so $[K:\Q]=n$ and $(1,\theta,\ldots,\theta^{n-1})$ is a basis. For $a\in K^*=K\setminus\{0\}$, define $M_a:K\to K$ by $M_a(v)=av$. The injection $a\mapsto M_a$ embeds $K^*$ into the group of invertible $\Q$-linear maps of $K$. This is the source's $\operatorname{Aut}_{\Q}(K)$, not the usually much smaller group of field automorphisms. Relative to the stated basis, write $G=\operatorname{GL}(n,\Q)$ and $T(d)=\{[M_a]:a\in K^*\}$, described in the source as an unsplit maximal torus.

The source considers subgroups $H\leq G$ containing $T(d)$ and some ``one-dimensional transformation'', says that they are rich in elementary transvections and so determine an associated net $\sigma(H)$, and asks whether
\[
H\leq N_G(E(\sigma(H))).
\]
Here, as supplied, $G(\sigma)$ is the net group and $E(\sigma)$ is generated by all transvections in $G(\sigma)$. A transvection in the standard linear sense is $I+v\lambda$, with nonzero vector $v$, nonzero linear functional $\lambda$ and $\lambda(v)=0$; elementary transvections include $I+a e_{ij}$ for $i\ne j$ and nonzero $a\in\Q$. For any subgroup $E\leq G$, its normalizer is
\[
N_G(E)=\{g\in G:gEg^{-1}=E\}.
\]
Thus the question is whether every $h\in H$ conjugates the associated transvection subgroup onto itself.

The notions of a rich subgroup, the construction $H\mapsto\sigma(H)$ and the definition of $G(\sigma)$ are not supplied in the record, so the quantified object $E(\sigma(H))$ cannot be fully reconstructed without the cited background. Also ``one-dimensional transformation'' may intend the standard condition $\operatorname{rank}(h-I)=1$, but that condition is not explicitly stated and is not silently substituted. The source does not explicitly delimit $n$; $n\geq2$ is declared above as the nontrivial extension range.

The background attributes the problem to V. A. Koibaev in Kourovka Notebook 21.77 (2026) and reports the inclusion already proved for $n=2$. The literature triage, checked 2026-08-17, labels the general question open and verifies no further resolution. Its generic suggestion about Zariski/arithmetic closure does not replace the actual normalizer-inclusion target, for which no such closure is specified. The primary registry status is ``open''.
\subsection{Short English statement}
For a rational linear group containing the multiplication copy of a radical field and a specified rank-one-type transformation, must it normalize the transvection group associated with its net? The source reports the degree-two case known, but leaves the rich-subgroup/net definitions to specialized references, so those conventions require source review.
\subsection{Sources}
\begin{itemize}
\item UnsolvedMath local registry, permanent ID 2586, source code catalogue entry, ``Kourovka Notebook Problem 21.77'', New Problems, 21st issue (2026), with V. A. Koibaev named in the background. Dataset attribution: Ulam AI and the UnsolvedMath Contributors; curation and metadata under CC BY 4.0.
\item V. A. Koibaev, St. Petersburg Mathematical Journal 21, no. 5 (2010), 731--742; V. A. Koibaev and A. V. Shilov, Journal of Mathematical Sciences (New York) 171, no. 3 (2010), 380--385. The background cites rich transvections and the associated net; titles and direct URLs are not supplied.
\item V. A. Koibaev, Doklady Mathematics 41, no. 3 (1990), 414--416. The background attributes the $n=2$ inclusion to this item; no title or direct URL is supplied.
\item Kourovka Notebook, 21st issue (2026), Problem 21.77, supplied tracker locator \url{https://www.kourovka-notebook.org/}. The local literature triage records its check as 2026-08-17.
\end{itemize}
\end{document}
